\begin{pblock}{HSR: elevation from the whole wetting curve}
A 10\,m pixel on the waterline is not wet or dry — it is
\term{partially flooded}, and its NDWI is the linear mixture of its wet and
dry fractions. As the tide rises that fraction follows the pixel's own
internal hypsometry, so its NDWI against tide height is a sigmoid whose
parameters are physical, not fitted conveniences:

\vspace{-2mm}
\begin{equation*}
  % \text, not \mathrm: with a sans text font and no matching math roman,
  % \mathrm falls back to a METAFONT bitmap — a Type 3 font that prints
  % jagged at a metre. \text picks up the document face instead.
  \text{NDWI}_p(t) \;=\; a_p \;+\; b_p\,
  \Phi\!\left(\frac{h_t-\mu_p}{\sigma_p}\right)
\end{equation*}

\defn{$\mu_p$ — median pixel elevation}{Read from the \emph{entire} curve
rather than from one threshold crossing, so every clear observation
contributes and the estimate is sub-decimetre.}
\defn{$\sigma_p$ — sub-pixel relief}{The roughness \emph{inside} the
10\,m cell: a product no threshold-crossing method can produce,
because it discards exactly the mixed pixels that carry it.}
\defn{Fit quality — a Cramér--Rao bound per pixel}{An uncertainty that is
derived rather than assumed, and that doubles as automatic quality control.}

The amplitude terms $(a_p,b_p)$ have a closed-form solution per candidate
$(\mu,\sigma)$: the fit is a cheap grid sweep, tractable over millions of
pixels. With MAREA, $h_t$ is each band's \term{corrected clock}.

\vspace{2mm}
\pfig[0.73\linewidth]{02_curvas_reales}
\figcap{Three real pixels, low to high on the flat. The sigmoid is not
imposed on the data — it is what mixed pixels do as the tide crosses
them, and its centre $\mu_p$ \emph{is} the elevation.}
\end{pblock}
